%!TEX program = lualatex
\documentclass[ngerman,headheight=70pt]{scrartcl}
\usepackage[ngerman]{babel}
\usepackage{blindtext}
\usepackage[utf8]{luainputenc}
\usepackage{scrlayer-scrpage}
\usepackage[T1]{fontenc}
\usepackage{fontspec}
\usepackage{multicol}
\usepackage{enumitem}
\usepackage{epigraph}
\usepackage[german=quotes]{csquotes}
\usepackage[bottom=2cm,footskip=8mm]{geometry}
\setmainfont{Source Sans Pro}
\MakeOuterQuote{"}

\setlength{\parindent}{0pt}
\setlength{\parskip}{6pt}

\pagestyle{plain.scrheadings}
\thispagestyle{scrheadings}
\lohead{\vspace*{-1.5cm}\includegraphics[scale=0.4]{../images/stupa_logo_rot.pdf}}
\rohead{\\Präsidium des\\Studierendenparlaments}

\lofoot[Universität Hamburg · Präsidium des Studierendenparlaments\\\hrule Von-Melle-Park 5 · D-20146 Hamburg]{Universität Hamburg · Präsidium des Studierendenparlaments\\\hrule Von-Melle-Park 5 · D-20146 Hamburg}
\cofoot[]{}
\cohead[\thepage]{}

\renewcommand{\thesubsection}{TOP \arabic{subsection}}
\renewcommand{\thesubsubsection}{\arabic{subsubsection}.}
\setcounter{subsection}{-1}

\begin{document}
    UHH · StuPa-Präsidium · Von-Melle-Park 5 · D-20146 Hamburg

    \section*{Beschluss des Studierendenparlamentes vom 9. Juni 2016}
    \subsection*{Kein Werben fürs Sterben!}

    Das Studierendenparlament der Uni Hamburg begrüßt, dass das
    Bundeswehr-Career Center nach studentischer Intervention im Vorfeld auf die
    Beteiligung an der Jobmesse 2016 der Uni Hamburg doch verzichtet hat. Die
    Universität ist kein Ort für Kriegswerbung und Soldatenrekrutierung. Denn
    Wissenschaft hat die Aufgabe, die Wahrheit für die menschenwürdige Entwicklung
    der Welt zu bestimmen und zur Wirklichkeit zu machen. Die Wahrheit stirbt im
    Krieg jedoch bekanntlich zu erst. Die Mitglieder der Uni Hamburg haben sich
    auch deswegen "Internationalisierung von Bildung und Wissenschaft für eine
    friedliche und menschenwürdige Welt" im Leitbild der Uni als Ziel ihrer
    Entwicklung gesetzt. Zahlreiche Unimitglieder arbeiten wissenschaftlich in
    verschiedenen Fachbereichen und Instituten an der zivilen Entwicklung der
    Gesellschaft. Dieses Friedensengagement ist der konsequente Ausdruck des
    Interesses der Mehrheit und die konsequente Fortführung der Erkenntnisse aus
    dem ersten und zweiten Weltkrieg: Nie wieder Faschismus! Nie wieder Krieg!

    Bis heute ist die Friedensbefürwortung mehrheitlich in der Bevölkerung
    verbreitet. Die Bundeswehr hat zunehmend mit dieser verbreiteten
    Kriegsablehnung zu kämpfen. So ist es erschwert, Geschäfte mit Krieg, Militär
    und Waffen zu machen und durchzusetzen, dass "im Notfall auch militärischer
    Einsatz notwendig ist, um unsere Interessen zu wahren, zum Beispiel freie
    Handelswege", wie es Horst Köhler, damaliger Bundespräsident, im Jahr 2010
    offen aussprach. Zu einer Berufs- und Angriffsarmee umstrukturiert, hat die
    Bundeswehr besonders Schwierigkeiten junge Leute für die Ausbildung anzuwerben.
    Deswegen investiert die Bundeswehr Millionen Euro öffentlicher Mittel in
    PR-Maßnahmen und versucht verstärkt in Schulen, Hochschulen und bei Jobmessen
    aufzutreten. Dabei versucht die Bundeswehr auch auszunutzen, dass staatliches
    Handeln mit neoliberaler Arbeitsmarktpolitik und dem Abbau sozialer
    Errungenschaften zu erheblicher und begründeter sozialer Verunsicherung
    beiträgt, indem sie sich demgegenüber als sicherer Arbeitgeber präsentiert.
    Gegen diesen Zynismus gibt es jedoch fast immer und meist erfolgreich
    Widerstand und Aufklärung aus der Bevölkerung, Friedensbewegung und den
    Gewerkschaften. Denn die Bundeswehr ist kein normaler Arbeitsgeber. Bei ihr
    sollen junge Menschen zu Soldaten und Soldatinnen, sprich zum Töten und damit
    zu potentiellen staatlich beauftragten Mördern, ausgebildet werden. Über 3000
    deutsche Soldaten und Soldatinnen befinden sich aktuell in 16 Auslandseinsätzen
    der Bundeswehr, vom Kosovo bis Somalia, von Afghanistan bis Mali.

    Die Bundeswehr hat an der Universität, die die wissenschaftliche Institution
    zur Wahrheitsfindung ist, nichts zu suchen. Das Studierendenparlament
    kritisiert das Stellenwerk der Uni Hamburg Marketing GmbH auf das Schärfste,
    der Bundeswehr auf der Jobmesse 2016 einen Stand genehmigt zu haben.
    Krieg ist die am meisten zugespitzte Form und der brutalste Ausdruck der
    Konkurrenzgesellschaft. Das Studierendenparlament fordert die Uni Hamburg
    Marketing GmbH auf, die kommenden Jobmessen an der Uni Hamburg so zu
    gestalten und auszurichten, dass den Studierenden ein sozialkritischer Zugang
    zum Arbeitsmarkt ermöglicht wird und die Bundeswehr nie wieder in Erwägung
    ziehen kann, an die Uni Hamburg für Rekrutierung von Soldaten und Soldatinnen
    zu kommen.

    Das Studierendenparlament nimmt die Auseinandersetzung mit der Bundeswehr im
    Vorfeld der Jobmesse 2016 zum Anlass erneut die Forderung nach einer
    Zivilklausel im Hamburgischen Hochschulgesetz zu bekräftigen. Mit einer
    Zivilklausel ist der verbindliche Maßstab zu setzen, dass Forschung, Studium
    und Lehre stets auf eine zivile Entwicklung gerichtet sein sollen. Sie ist
    damit auch Schutz der Wissenschaften vor der Macht des
    militärisch-industriellen Komplexes. Das Studierendenparlament ruft alle und
    insbesondere den AStA auf, das Engagement für die weltweite
    Friedensentwicklung zu stärken. Der AStA wird beauftragt, im Bündnis
    "Bildung ohne Bundeswehr" mitzuarbeiten.

    Das Studierendenparlament ruft dazu auf, an der Kundgebung "Kein Tag der Bundeswehr",
    am 11. Juni 2016 von 10 bis 18 Uhr, teilzunehmen.

\end{document}
